\chapter{Introduction}
\label{ch:introduction}

The rapid proliferation of Artificial Intelligence (AI) has ushered in a new era of technological capability, with systems now deployed in high-stakes domains ranging from criminal justice and healthcare to finance and autonomous transportation. While the potential benefits are immense, the opacity, complexity, and scale of these systems introduce significant ethical risks. Incidents of algorithmic bias, discriminatory outcomes, privacy violations, and lack of accountability have become increasingly common, eroding public trust and highlighting the urgent need for systematic governance and oversight~\cite{zhang2023fairness,jobin2019landscape}. In response, governments, academic institutions, and industry bodies have developed a variety of ethical principles and regulatory frameworks. However, a significant gap persists between high-level principles and their practical implementation in the software development lifecycle~\cite{mittelstadt2019principles}. Developers and organizations often lack the specific tools and methodologies to translate abstract ethical goals into concrete, measurable, and auditable engineering requirements.

This Final Year Project addresses this critical gap by designing, implementing, and evaluating the Multi-stakeholder Ethical Decision Framework (MEDF). MEDF is a decision-support platform that provides a structured and reproducible workflow for evaluating the ethical dimensions of AI applications. It operationalizes abstract principles from leading governance frameworks, the EU AI Act Assessment List for Trustworthy AI (ALTAI)~\cite{aihleg2020altai}, the NIST AI Risk Management Framework (RMF)~\cite{nist2023rmf}, and Singapore's Model AI Governance Framework (MGAF)~\cite{pdpc2020mgaf}, into a unified set of six quantifiable dimensions: transparency and explainability, fairness and non-discrimination, safety and robustness, privacy and data governance, human agency and oversight, and accountability.

This report details the complete lifecycle of the MEDF project, from initial problem definition to final implementation and evaluation. It begins by reviewing the background and motivation for ethical AI governance, followed by a comprehensive literature review of existing frameworks and decision-making models. The core of the report is dedicated to the design and architecture of the MEDF system, including its three-tier structure, the algorithmic foundations of its scoring and conflict-resolution engines (TOPSIS, WSM, NSGA-II), and the rationale behind critical engineering decisions. The practical utility of the framework is demonstrated through the in-depth analysis of three real-world case studies: the Metropolitan Police's use of live facial recognition, iTutorGroup's automated hiring system, and the deployment of DeepMind's Streams application in the UK's National Health Service (NHS). The report concludes with a discussion of the project's limitations, its contributions to the field, and potential avenues for future research. The final output is a completely functional, open-source software artifact that serves as both a practical tool for AI governance and a reproducible research platform for future work in computational ethics.
